En un medio de volumen constante, $A+B\xrightarrow{k_1}C\xrightarrow{k_2}D$, con cinética de acción de masas, $k_1,k_2>0$. Inicialmente $A(0)=A_0$, $B(0)=B_0$, $C(0)=D(0)=0$ y $A_0\gg B_0>0$. Para las expresiones explícitas use la aproximación de reactivo en exceso $k_1A(t)\simeq k_1A_0=\lambda$, distinguiéndola de los balances de materia. Las cantidades se expresan como concentraciones compatibles con la estequiometría uno a uno.

\begin{enumerate}[label=(\alph*),ref=\alph*,leftmargin=*]
\item Dé las ecuaciones diferenciales y los balances de materia del sistema.
\item Determine las cantidades de los cuatro compuestos con el tiempo en esa aproximación.
\item Pruebe que, si $A_0=k_2/k_1$, el máximo de $C$ en el modelo reducido no depende de las constantes de reacción. Halle ese máximo y el instante en que se alcanza.
\end{enumerate}
